% Preface of the monograph (input by AG_monograph.tex)
\chapter*{Preface}
\addcontentsline{toc}{chapter}{Preface}
\markboth{Preface}{Preface}

\paragraph{What this book is about.} Sequencing does not read a genome; it reads its words. From
short reads one learns, at best, the multiset of $k$-mers of the genome and of its reverse
complement, and everything an assembler, a variant caller or a sequence encoder does is constrained
by what that multiset determines. This book is about the exact size and shape of that constraint.
Given the data, the set of genomes compatible with them is a \emph{fibre}. We count fibres exactly,
describe the moves that connect their members, prove where exact counting stops being tractable, and
show where the same fibres reappear in modern machine-learning representations of biological
sequences.

\paragraph{One strand (Part I).} For a single strand the fibre is the set of Eulerian circuits of a
de Bruijn multigraph, and its size is a determinant, by the theorem of van Aardenne-Ehrenfest,
de Bruijn, Smith and Tutte. The determinant factors into a local part, fixed by the branch
statistics, and a global part: the order of the graph's sandpile group. Chapter~1 shows that the
global part is not determined by the branch statistics. Chapter~2 turns the group into a coordinate
system on reconstructions, one chip per rearrangement. Chapter~3 splits it along the repeats into a
part that longer words remove and a part that no word length removes. Chapter~4 shows where pairwise
repeat data stop describing it.

\paragraph{Two strands (Part II).} Real data do not know which strand a word came from. The
reverse-complement double cover keeps every single-strand lemma (Chapter~5). But its involution is an
anti-automorphism of the Laplacian (Chapter~6), and the double-stranded count is no longer one
determinant but a sum of determinants over the lattice points of a box (Chapter~7). For the phage
$\phi$X174 at word length $10$ the box has $2\,481\,687\,900$ points (Chapter~8), and the count is
$31\,925\,753\,246\,212$ (Chapter~9). Chapter~10 proves that this is not a failure of ingenuity:
the double-stranded count is $\#$P-hard, so no closed formula computes it unless $\mathrm{FP}=\#\mathrm P$.

\paragraph{Fibres and moves (Part III).} The members of a fibre are connected by explicit moves:
transpositions between direct repeats and inversions between inverted repeats. Chapter~11 computes
the lattice that the inversions span, by a Tate-cohomology dichotomy, and states a conjecture on
generation. Chapter~12, written with Zhengyi Chen, shows that the fibres are exactly what local
sequence encoders cannot see. It computes them for the whole human proteome and constructs explicit
proteins that a widely used composition encoder cannot tell apart. Chapters~13 and~14 treat pooled
data: polyploid phasing as a decomposition of one flow, and pan-genomes through their recombination
group.

\paragraph{Algorithms (Part IV).} Hardness is a statement about the worst case. Real genomes are
tractable when their tangled parts are small, and the book names the reason each time. Chapter~15
reduces the $4.6$-million-vertex graph of \emph{E.~coli} to $115$ vertices by contracting everything
except repeats. Chapter~16 computes the single-strand determinant in linear time along the snarl tree
of a genome graph. Chapter~17 describes the polytope of isoform abundances compatible with splicing
data.

\paragraph{How the results are checked.} Every chapter comes with a verification script, and every
number in the text, the tables and the figures is printed by one of them. The scripts check three
kinds of statement, and the book keeps them apart:
\begin{itemize}
\item \emph{theorems}, which are proved in the text; the scripts test them on exhaustive or random
instances, and a failure would indicate an error in the proof or in the code;
\item \emph{computations}, such as $\#\DS(\phi\mathrm X174,10)$, which the scripts produce and
cross-check by independent methods;
\item \emph{conjectures and empirical statements}, which are labelled as such and for which the
scripts report the evidence and nothing more.
\end{itemize}
The repository \url{https://github.com/Ruqing1963/algebraic-genomics}, archived under
DOI~\href{https://doi.org/\AGdoi}{\texttt{\AGdoi}}, contains the papers, the scripts, the input data
and the reference transcripts. The script \texttt{build\_all.py} compiles every chapter and reruns
every check; \texttt{-{}-quick} takes about half an hour, \texttt{-{}-full} an afternoon. The state of
that build for this release is recorded at the end of the book.

\paragraph{What the book corrects.} Several chapters began as specifications that turned out to be
wrong in part: a rank that was not the number of inverted repeats, a normalisation that was missing,
a tensor network that was a Schur complement, a polynomial formula that cannot exist. Each chapter
says which parts of its specification failed and why. We have kept these corrections in the text
because they are part of the result.

\paragraph{How to read it.} Part~I is self-contained. Part~II assumes Chapters~1 and~3. Part~III can be
read after Chapter~7. Part~IV uses Chapter~3 (Chapter~15) and the BEST theorem (Chapters~16, 17). The
chapters were written as papers; we have unified the notation, numbering and bibliography, in Part~I replaced repeated definitions by references to Chapter~1, and each
chapter still opens with a summary that can be read alone.

\paragraph{Acknowledgements.} Chapter~12 is joint work with Zhengyi Chen (Guangxi Key Laboratory of
Drug Discovery and Optimization, School of Pharmacy, Guilin Medical University), and is included with
the co-author's agreement. The phage, bacterial and human data are those of NCBI RefSeq, GenBank and UniProtKB.

\bigskip
\noindent Ruqing Chen\hfill Rivi\`ere-Beaudette, Qu\'ebec, \AGrelease
